% GENERATED FILE. Edit canonical lessons, then run npm run generate:books.
% canonical-source-hash: fnv1a64:314f5e8a5ad35dcf
% canonical-lessons: ML-C122-really, ML-C122-nothing, ML-C122-ohno, ML-C122-thesame, ML-C122-ohisee

\chapter{Quick Replies}
\label{ch:ml-quick-replies}
\hlchaptermodality{122}

\begin{chapteropening}
\textbf{By the end of this chapter:} \emph{I can say really?, nothing, oh no!, the same, and oh I see.}

Complete the last lesson of chapter 122: i can say really?, nothing, oh no!, the same, and oh I see.
\end{chapteropening}

\section[śarikkuṁ]{\ml{ശരിക്കും} (śarikkuṁ) — really}

\label{lesson:ML-C122-really}



\begin{quote}
Before the new one: say the Malayalam for a station, then the Malayalam for a map.
\end{quote}

\subsection*{You'll want to know: \ml{ശരിക്കും}}
\textbf{\ml{ശരിക്കും}} (\emph{śarikkuṁ}) — \textquotedblleft{}really, truly\textquotedblright{}, built on \textbf{\ml{ശരി}} (\emph{śari}), \textquotedblleft{}right, okay\textquotedblright{}. \textbf{\ml{ശരിക്കും}?} — \textquotedblleft{}really?\textquotedblright{}.

\begin{grammarlens}[title={What you've built}]
Okay, turned into really.
\end{grammarlens}

\subsection*{Your turn}
\begin{itemize}
\raggedright
  \item \emph{Say it:} \emph{śarikkuṁ}
  \item \emph{Say it:} \emph{śarikkuṁ}, once more
  \item \emph{Say it:} say \emph{śarikkuṁ?}
  \item \emph{Recall:} say the Malayalam for a station, then the Malayalam for a map, then say \emph{śarikkuṁ} again
\end{itemize}

\begin{tcolorbox}[breakable,colback=teal!4,colframe=teal!35!black,title={Before you move on}]
What does \ml{ശരിക്കും} mean? (\textquotedblleft{}Really\textquotedblright{}.) Say it once more.
\end{tcolorbox}
\section[onnumilla]{\ml{ഒന്നുമില്ല} (onnumilla) — nothing}

\label{lesson:ML-C122-nothing}



\begin{quote}
Before the new one: say the Malayalam for a map, then the Malayalam for really.
\end{quote}

\subsection*{You'll want to know: \ml{ഒന്നുമില്ല}}
\textbf{\ml{ഒന്നുമില്ല}} (\emph{onnumilla}) — \textquotedblleft{}nothing\textquotedblright{}, literally \textquotedblleft{}not even one\textquotedblright{}: \textbf{\ml{ഒന്ന്}}, \textquotedblleft{}one\textquotedblright{}, \textbf{‑\ml{ഉം}}, \textquotedblleft{}even\textquotedblright{}, \textbf{\ml{ഇല്ല}}, \textquotedblleft{}not\textquotedblright{}.

\begin{grammarlens}[title={What you've built}]
One, even, not.
\end{grammarlens}

\subsection*{Your turn}
\begin{itemize}
\raggedright
  \item \emph{Say it:} \emph{onnumilla}
  \item \emph{Say it:} \emph{onnumilla}, once more
  \item \emph{Say it:} answer \emph{onnumilla}
  \item \emph{Recall:} say the Malayalam for a map, then the Malayalam for really, then say \emph{onnumilla} again
\end{itemize}

\begin{tcolorbox}[breakable,colback=teal!4,colframe=teal!35!black,title={Before you move on}]
What does \ml{ഒന്നുമില്ല} mean? (\textquotedblleft{}Nothing\textquotedblright{}.) Say it once more.
\end{tcolorbox}
\section[ayyō]{\ml{അയ്യോ} (ayyō) — oh no!}

\label{lesson:ML-C122-ohno}



\begin{quote}
Before the new one: say the Malayalam for really, then the Malayalam for nothing.
\end{quote}

\subsection*{You'll want to know: \ml{അയ്യോ}}
\textbf{\ml{അയ്യോ}} (\emph{ayyō}) — \textquotedblleft{}oh no!\textquotedblright{}, \textquotedblleft{}oh dear!\textquotedblright{}, heard everywhere in the south when something goes wrong.

\begin{grammarlens}[title={What you've built}]
The sound of a small disaster.
\end{grammarlens}

\subsection*{Your turn}
\begin{itemize}
\raggedright
  \item \emph{Say it:} \emph{ayyō}
  \item \emph{Say it:} \emph{ayyō}, once more
  \item \emph{Say it:} say \emph{ayyō!}
  \item \emph{Recall:} say the Malayalam for really, then the Malayalam for nothing, then say \emph{ayyō} again
\end{itemize}

\begin{tcolorbox}[breakable,colback=teal!4,colframe=teal!35!black,title={Before you move on}]
What does \ml{അയ്യോ} mean? (\textquotedblleft{}Oh no!\textquotedblright{}.) Say it once more.
\end{tcolorbox}
\section[atutanne]{\ml{അതുതന്നെ} (atutanne) — that very one, the same}

\label{lesson:ML-C122-thesame}



\begin{quote}
Before the new one: say the Malayalam for nothing, then the Malayalam for oh no.
\end{quote}

\subsection*{You'll want to know: \ml{അതുതന്നെ}}
\textbf{\ml{അതുതന്നെ}} (\emph{atutanne}) — \textquotedblleft{}that very one, exactly that\textquotedblright{}: \textbf{\ml{അത്}} (\emph{atŭ}), \textquotedblleft{}that\textquotedblright{}, with \textbf{\ml{തന്നെ}}, \textquotedblleft{}itself\textquotedblright{}.

\begin{grammarlens}[title={What you've built}]
That, made emphatic.
\end{grammarlens}

\subsection*{Your turn}
\begin{itemize}
\raggedright
  \item \emph{Say it:} \emph{atutanne}
  \item \emph{Say it:} \emph{atutanne}, once more
  \item \emph{Say it:} say \emph{atutanne!}
  \item \emph{Recall:} say the Malayalam for nothing, then the Malayalam for oh no, then say \emph{atutanne} again
\end{itemize}

\begin{tcolorbox}[breakable,colback=teal!4,colframe=teal!35!black,title={Before you move on}]
What does \ml{അതുതന്നെ} mean? (\textquotedblleft{}That very one, the same\textquotedblright{}.) Say it once more.
\end{tcolorbox}
\section[ōhō]{\ml{ഓഹോ} (ōhō) — oh, I see}

\label{lesson:ML-C122-ohisee}



\begin{quote}
Before the new one: say the Malayalam for oh no, then the Malayalam for the same.
\end{quote}

\subsection*{You'll want to know: \ml{ഓഹോ}}
\textbf{\ml{ഓഹോ}} (\emph{ōhō}) — \textquotedblleft{}oh, I see!\textquotedblright{}, the sound of understanding arriving.

\begin{grammarlens}[title={What you've built}]
The happy opposite of ayyō.
\end{grammarlens}

\subsection*{Your turn}
\begin{itemize}
\raggedright
  \item \emph{Say it:} \emph{ōhō}
  \item \emph{Say it:} \emph{ōhō}, once more
  \item \emph{Say it:} say \emph{ōhō!}
  \item \emph{Recall:} say the Malayalam for oh no, then the Malayalam for the same, then say \emph{ōhō} again
\end{itemize}

\begin{tcolorbox}[breakable,colback=teal!4,colframe=teal!35!black,title={Before you move on}]
What does \ml{ഓഹോ} mean? (\textquotedblleft{}Oh, I see\textquotedblright{}.) Say it once more.
\end{tcolorbox}
